\documentclass[11pt]{article}
\usepackage[margin=1in]{geometry}
\usepackage{amsmath,amssymb,amsthm}
\usepackage{xurl}
\usepackage{hyperref}
\sloppy
\let\oldus\_ \renewcommand{\_}{\oldus\allowbreak}
\newtheorem{theorem}{Theorem}
\newtheorem{lemma}[theorem]{Lemma}
\theoremstyle{definition}
\newtheorem{definition}[theorem]{Definition}
\title{Conjecture 00000002805 is false:\\ the Schilder rate does not vanish on the Cameron--Martin unit ball}
\author{Jackmeson1}
\date{2026-10-05}
\begin{document}
\maketitle

\section{The conjecture}
Conjecture 00000002805 (English text, verbatim): ``Definition: Schilder's theorem: the LDP on
the path space of Brownian motion. Conjecture: Cameron--Martin rigidity of the Schilder rate: the
rate function vanishes if and only if the path lies in the Cameron--Martin unit ball, and the
spectral distribution of the rate on the sphere is an explicit chi-square-type law.'' The Chinese
text says the same.

\paragraph{Reading.} The conjecture is a conjunction. We refute its first clause, the
``Cameron--Martin rigidity'' named in the title:
\[ I(\omega)=0 \iff \omega \text{ lies in the Cameron--Martin unit ball}, \tag{$\ast$}\]
where $I$ is the rate function of Schilder's theorem. We refute $(\ast)$ for the closed unit ball
$\{\|\omega\|_H\le 1\}$ and for the open unit ball $\{\|\omega\|_H<1\}$, in every dimension
$d\ge 1$ and for every time horizon $T>0$. The second clause (the ``spectral distribution on the
sphere'') is not interpreted and not used; refuting one conjunct refutes the conjunction.

\section{Definitions (Schilder's theorem)}
We use the statement of Schilder's theorem as given on Wikipedia (``Schilder's theorem'', retrieved
2026-10-05). Let $B$ be a standard Brownian motion in $\mathbb R^d$ started at $0$, and
$C_0=C_0([0,T];\mathbb R^d)$ the continuous paths $\omega$ on $[0,T]$ with $\omega(0)=0$. For
$\varepsilon>0$ the laws of $\sqrt\varepsilon B$ satisfy the large deviation principle on $C_0$
``with good rate function $I : C_0 \to \mathbb R\cup\{+\infty\}$ given by
$I(\omega)=\frac12\int_0^T|\dot\omega(t)|^2\,dt$ if $\omega$ is absolutely continuous, and
$I(\omega)=+\infty$ otherwise''. The Cameron--Martin space is ``the subspace of absolutely
continuous functions whose derivative is in $L^2$''. Its norm is
$\|\omega\|_H=\big(\int_0^T|\dot\omega(t)|^2\,dt\big)^{1/2}$, so $I(\omega)=\frac12\|\omega\|_H^2$
on $H$.

\begin{definition}[as in Lean]
$\mathbb R^d$ is \texttt{EuclideanSpace R (Fin d)}; paths are functions $\omega:\mathbb R\to\mathbb R^d$
looked at on $[0,T]$.
\begin{itemize}
\item \texttt{PathSpace d T}: $\omega$ continuous on $[0,T]$ and $\omega(0)=0$.
\item \texttt{CameronMartin d T}: $\omega(0)=0$, $\omega$ absolutely continuous on $[0,T]$
  (Mathlib's \texttt{AbsolutelyContinuousOnInterval}), and $\dot\omega=$ \texttt{deriv}~$\omega$ in
  $L^2([0,T];\mathbb R^d)$ (\texttt{MemLp (deriv omega) 2}).
\item \texttt{cmNormSq T omega} $=\int_{[0,T]}\|\dot\omega(t)\|^2\,dt\in[0,\infty]$ (lower Lebesgue
  integral), i.e.\ $\|\omega\|_H^2$.
\item \texttt{cmClosedUnitBall}, \texttt{cmOpenUnitBall}: $\omega\in H$ with
  $\|\omega\|_H^2\le 1$, resp.\ $<1$.
\item \texttt{schilderRate d T omega} $=\frac12\,\|\omega\|_H^2$ if $\omega$ is absolutely continuous on
  $[0,T]$, and $+\infty$ otherwise (values in $[0,\infty]$).
\end{itemize}
\end{definition}
Using the lower integral makes $I(\omega)=+\infty$ automatically when $\omega$ is absolutely
continuous but $\dot\omega\notin L^2$, which agrees with the textbook rate.

\section{The counterexample}
Let $e_0$ be the first standard basis vector of $\mathbb R^d$ ($d\ge1$), put
$w=\frac{1}{2\sqrt T}e_0$ and $\omega_0(t)=t\,w$.

\begin{theorem}
$\omega_0\in C_0$, $\omega_0\in H$, $\|\omega_0\|_H^2=\frac14$ and $I(\omega_0)=\frac18$.
Hence $\omega_0$ lies in the open (and so in the closed) Cameron--Martin unit ball but
$I(\omega_0)\ne0$, and $(\ast)$ fails for both balls.
\end{theorem}
\begin{proof}
$\omega_0$ is linear, so it is continuous, $\omega_0(0)=0$, it is $C^1$ and hence absolutely
continuous on $[0,T]$, and $\dot\omega_0(t)=w$ for every $t$. A constant is in $L^2$ of the finite
measure space $[0,T]$, so $\omega_0\in H$. Since $\|e_0\|=1$, $\|w\|=\frac{1}{2\sqrt T}$ and
\[\|\omega_0\|_H^2=\int_0^T\frac{1}{4T}\,dt=\frac14<1,\qquad I(\omega_0)=\frac12\cdot\frac14=\frac18\neq0.\]
If $(\ast)$ held for all $\omega\in C_0$, the implication ``in the ball $\Rightarrow I=0$'' applied
to $\omega_0$ would give $I(\omega_0)=0$, a contradiction.
\end{proof}

\paragraph{Remark (what is true).} $I(\omega)=\frac12\|\omega\|_H^2$ vanishes exactly at the zero
path, while the Cameron--Martin unit ball is the sublevel set $\{I\le\frac12\}$. This remark is not
formalized and is not used.

\section{Lean formalization}
File \texttt{lean/Conjecture2805/Basic.lean} (148 lines, only import \texttt{Mathlib}), namespace
\texttt{C2805}. Objects: \texttt{PathSpace}, \texttt{CameronMartin}, \texttt{cmNormSq},
\texttt{cmClosedUnitBall}, \texttt{cmOpenUnitBall}, \texttt{schilderRate} as above;
\texttt{dir d T} $=\frac{1}{2\sqrt T}e_0$; \texttt{straightPath d T} $=t\mapsto t\cdot$\texttt{dir d T}.
Lemmas: \texttt{straightPath\_mem\_pathSpace}, \texttt{straightPath\_mem\_cameronMartin}
(absolute continuity from \texttt{ContDiffOn.absolutelyContinuousOnInterval}),
\texttt{cmNormSq\_straightPath} ($=1/4$), \texttt{schilderRate\_straightPath} ($=1/8$),
\texttt{straightPath\_mem\_open\_ball}, \texttt{straightPath\_mem\_closed\_ball},
\texttt{schilderRate\_straightPath\_ne\_zero}. Main theorems, for every $d$ with
\texttt{[NeZero d]} and every $T>0$:
\begin{itemize}
\item \texttt{schilder\_rigidity\_false}: \texttt{not (forall omega in PathSpace d T,
  (schilderRate d T omega = 0 <-> omega in cmClosedUnitBall d T))};
\item \texttt{schilder\_rigidity\_false\_open}: the same with \texttt{cmOpenUnitBall}.
\end{itemize}
\texttt{\#print axioms} for both main theorems and for \texttt{schilderRate\_straightPath} gives
\texttt{[propext, Classical.choice, Quot.sound]}; no \texttt{sorry}, no \texttt{native\_decide}.
Checked with Lean 4.33.1 and Mathlib v4.33.1 (\texttt{lake build}).

\section{Relation to the earlier submission}
An earlier submission (PR \#277, orionsheep) gave the same counterexample $\omega(t)=t/2$ but was
closed. The reviewer wrote that its main theorem ``is entirely vacuous $\mathbb N$ trivia
($\forall c>0$, $0<c\cdot c$, $1\le2$, $1\ne0$, $2=1\cdot2$, $1\ne2$); the Schilder rate,
Cameron--Martin space, sphere, and chi-square law never appear in Lean''. This submission defines
in Lean the path space $C_0([0,T];\mathbb R^d)$, the Cameron--Martin space (absolutely continuous,
derivative in $L^2$), the Cameron--Martin norm and both unit balls, and Schilder's rate function
with the value $+\infty$ off the absolutely continuous paths; it proves the counterexample for
every $d\ge1$ and $T>0$. Unlike PR \#277 it makes no claim about the second clause (the sphere and
the chi-square law), which is left uninterpreted.

\section*{Source}
Wikipedia, ``Schilder's theorem'', \url{https://en.wikipedia.org/wiki/Schilder%27s_theorem}
(definitions quoted above).
\end{document}
